% This is the ADASS_template.tex LaTeX file, 25th Oct 2018.
% It is based on the ASP general author template file, but modified to reflect the specific
% requirements of the ADASS proceedings.
% Copyright 2014, Astronomical Society of the Pacific Conference Series
% Revision:  14 August 2014

% To compile, at the command line positioned at this folder, type:
% latex ADASS_template
% latex ADASS_template
% dvipdfm ADASS_template
% This will create a file called ADASS_template.pdf

\documentclass[11pt,twoside]{article}

% Do NOT use ANY packages other than asp2014. 
\usepackage{asp2014}

\aspSuppressVolSlug
\resetcounters

% References must all use BibTeX entries in a .bibfile.
% References must be cited in the text using \citet{} or \citep{}.
% Do not use \cite{}.
% See ManuscriptInstructions.pdf for more details
\bibliographystyle{asp2014}

% The ``markboth'' line sets up the running heads for the paper.
% 1 author: "Surname"
% 2 authors: "Surname1 and Surname2"
% 3 authors: "Surname1, Surname2, and Surname3"
% >3 authors: "Surname1 et al."
% Replace ``Short Title'' with the actual paper title, shortened if necessary.
% Use mixed case type for the shortened title
% Ensure shortened title does not cause an overfull hbox LaTeX error
% See ASPmanual2010.pdf 2.1.4  and ManuscriptInstructions.pdf for more details
\markboth{A. Luise, M. Lombardi, N. Parmiggiani}{Explaibale XRF Element Detection}

\begin{document}

\title{Explainable XRF Element Detection with Embedding-Space k-NN and Physics-Based Masks for the HARLOCK Lunar Rover}

% Note the position of the comma between the author name and the 
% affiliation number.
% Authors surnames should come after first names or initials, eg John Smith, or J. Smith.
% Author names should be separated by commas.
% The final author should be preceded by "and".
% Affiliations should not be repeated across multiple \affil commands. If several
% authors share an affiliation this should be in a single \affil which can then
% be referenced for several author names.
% See ManuscriptInstructions.pdf and ASPmanual2010.pdf 3.1.4 for more details

\author{Alberto Luise,$^1$ Michele Lombardi,$^1$ Nicolò Parmiggiani,$^2$ Andrea Bulgarelli,$^2$ Samuele Salti,$^1$ Sara Trevisan,$^3$ Federico Dogo,$^3$ Luca Castaldini,$^2$ Ambra Di Piano,$^2$ Francesca Esposito,$^4$ Antonio Femia,$^2$ Fabrizio Fiore,$^3$ Eduardo Medinaceli,$^2$ Ghulam Mustafa,$^1$ Alessio Trois$^5$ \\
\affil{$^1$University of Bologna, Bologna, BO, Italy}
\affil{$^2$INAF / OAS, Bologna, BO, Italy}
\affil{$^3$INAF/OA Trieste, Trieste, TS, Italy}
\affil{$^4$INAF/OA Napoli, Napoli, NA, Italy}
\affil{$^5$INAF/OA Cagliari, Cagliari, CA, Italy}}

% This section is for ADS Processing.  There must be one line per author.

\paperauthor{Alberto Luise}{alberto.luise@unibo.it}{}{University of Bologna}{DISI}{Bologna}{BO}{40136}{Italy}
\paperauthor{Michele Lombardi}{michele.lombardi2@unibo.it}{}{University of Bologna}{DISI}{Bologna}{BO}{40136}{Italy}
\paperauthor{Nicolò Parmiggiani}{nicolo.parmiggiani@inaf.it}{}{INAF / OAS}{ }{Bologna}{BO}{}{Italy}
\paperauthor{Andrea Bulgarelli}{andrea.bulgarelli@inaf.it}{}{INAF / OAS}{ }{Bologna}{BO}{}{Italy}
\paperauthor{Samuele Salti}{samuele.salti@unibo.it}{}{University of Bologna}{DISI}{Bologna}{BO}{40136}{Italy}
\paperauthor{Sara Trevisan}{sara.trevisan@unitn.it}{}{INAF/OA Trieste}{ }{Trieste}{TS}{}{Italy}
\paperauthor{Federico Dogo}{federico.dogo@inaf.it}{}{INAF/OA Trieste}{ }{Trieste}{TS}{}{Italy}
\paperauthor{Luca Castaldini}{luca.castaldini@inaf.it}{}{INAF / OAS}{ }{Bologna}{BO}{}{Italy}
\paperauthor{Ambra Di Piano}{}{}{INAF / OAS}{ }{Bologna}{BO}{}{Italy}
\paperauthor{Francesca Esposito}{francesca.esposito@inaf.it}{}{INAF/OA Napoli}{ }{Napoli}{NA}{}{Italy}
\paperauthor{Antonio Femia}{antonio.femia@inaf.it}{}{INAF / OAS}{ }{Bologna}{BO}{}{Italy}
\paperauthor{Fabrizio Fiore}{fabrizio.fiore@inaf.it}{}{INAF/OA Trieste}{ }{Trieste}{TS}{}{Italy}
\paperauthor{Eduardo Medinaceli}{eduardo.medinaceli@inaf.it}{}{INAF / OAS}{ }{Bologna}{BO}{}{Italy}
\paperauthor{Ghulam Mustafa}{Ghulam Mustafa}{}{University of Bologna}{DISI}{Bologna}{BO}{40136}{Italy}
\paperauthor{Alessio Trois}{alessio.trois@inaf.it}{}{INAF/OA Cagliari}{ }{Cagliari}{CA}{}{Italy}


% There should be one \aindex line (commented out) for each author. These are used to
% build up the author index for the Proceedings. The surname must come first, followed by
% initials. Note the use of ~ before each initial to control spacing.
% The \author entries (see above) have surname last. These \aindex entries have
% surname first.

%\aindex{Luise,~A.}
%\aindex{Lombardi,~M.}
%\aindex{Parmiggiani,~N.}
%\aindex{Bulgarelli,~A.}
%\aindex{Salti,~S.}
%\aindex{Trevisan,~S.}
%\aindex{Dogo,~F.}
%\aindex{Castaldini,~L.}
%\aindex{Di Piano,~A.}
%\aindex{Esposito,~F.}
%\aindex{Femia,~A.}
%\aindex{Fiore,~F.}
%\aindex{Medinaceli,~E.}
%\aindex{Mustafa,~G.}
%\aindex{Trois,~A.}

\begin{abstract}
%In this work we propose a framework that converts any Convolutional spectral regressor into an explainable predictor without hindering its performance on multiple metrics. Deep neural networks are often considered "black boxes", and the lack of transparency can limit their adoption in scientific contexts; this is one of the challenges in the development of RHEA, the Rare-earth \& Heavy Elements Analyzer of the HARLOCK project within the PRORIS program. It combines gamma-ray spectroscopy, multispectral imaging, ultraviolet fluorescence screening and X-ray fluorescence analysis to identify Rare-Earth Elements (REE) deposits and associated elemental and mineralogical signatures across different spatial scales. 
%Our approach exploits the CNN's latent representation and preservation of spatial features: after training, the final regression layer of the backbone model is replaced by a K-Nearest Neighbors (k-NN) regressor operating directly on the learned embedding space; as a result, every prediction is accompanied by an explanation in the form of the most similar spectra observed during training. This works particularly well thanks to a custom loss function designed to optimize the quality of the learned embedding space, without hindering performances. Furthermore, in our model, an attention mask is learned for each predicted element starting from a physics-based initialization, providing easily readable heatmaps as well in the form of a mask-by-weights product. Experimental results show that, in a regression task estimating the composition of soil samples, the average R2 score over the predicted elements stays around 0.9 for both the original approach and the explainable Deep k-NN model.
Deep neural networks are often considered "black boxes", and the lack of transparency can limit their adoption in scientific contexts; in this work, we propose a framework that converts any Convolutional regressor for mono-dimensional signals into an explainable predictor without hindering its performance. This is a contribution for RHEA, the Rare-earth \& Heavy Elements Analyzer of the HARLOCK rover project within the PRORIS program. Its payload leverages X-Ray Flourescence analysis (alongside other techniques) to identify Rare-Earth Elements (REE) deposits and associated mineralogical signatures across different spatial scales. 
Our approach exploits the CNN's latent representation and preservation of spatial features: after training, the final regression layer of the backbone model is replaced by a K-Nearest Neighbors (k-NN) regressor operating directly on the learned embedding space; as a result, every prediction is accompanied by an explanation in the form of the most similar spectra observed during training. This works particularly well thanks to a custom loss function designed to optimize the quality of the learned embedding space, without hindering performances. Furthermore, an attention mask is learned for each predicted element starting from a physics-based initialization, providing easily readable heatmaps as well in the form of a mask-by-weights product. Experimental results show that, in a regression task estimating the composition of soil samples, the average $R^2$ score over the predicted elements stays around 0.9 for both the original approach and the explainable Deep k-NN model.

\end{abstract}


\section{Introduction}
The HARLOCK project is located within the PROSIS program (\textit{www.proris.it}), coordinated by the Italian National Institute for Astrophysics (INAF) and the National Research Council (CNR). It aims to develop a  proof-of-concept for a fully autonomous lunar rover designed to scout and prospect lunar resources such as water, REEs and rare metals.
A key component of this mission is RHEA (\textit{Rare-earth \& Heavy Elements Analyzer}), the rover's scientific payload. Among the multiple techniques used by RHEA we want to focus \textbf{X-ray Fluorescence (XRF) analysis}, a non-destructive analytical technique used to determine the elemental composition of a material. It works by bombarding the sample with high-energy X-rays and measuring the energy and intensity of "secondary" X-rays, leveraging how each elements emits photons at a different rate and energy level.
This paper presents an explainable deep learning framework for spectral regression, aiming to determine the exact elemental composition of lunar samples from an XRF spectrum. We propose physics-injected attention masks, that stabilize training and provide interpretable feature attributions aligned with known atomic emission lines, and we implement a k-Nearest Neighbors (k-NN) regressor operating on a well-ordered latent embedding space replacing the last layer of a Convolutional Neural Network (CNN), yielding both accurate predictions and instance-level explanations from training data.

\section{Dataset Simulation}
Despite XRF analysis being widely adopted for a variety of tasks, 
%like \citet{PAINTINGS},
a dataset of lunar soil compositions bombarded with X-Rays is not readily available. For this reason, we developed a simulation pipeline using data from two of the few lunar samples already available in literature (\citet{Meyer2011LunarSample10072} and \citet{Meyer2011LunarSample14003}). A synthetic Lunar Soil dataset was generated by sampling elemental concentrations from a Dirichlet distribution centered around the two samples, aiming to to create a realistic enough environment; concentrations can range from 40\% (which is always Oxygen, given all the oxides present) down to the parts per billion, which aren't detectable. For this reason, we usually set a minimum threshold of $10^{-4}$.
Once the data was ready, it was used to generate the secondary X-ray spectra via the \textbf{Fundamental Parameters method}.

The pipeline models the physical XRF response of a geological sample in a very specific setup: utilizing the Fundamental Parameters method, the simulator estimates primary and secondary X-ray fluorescence based on i) the elemental and mineralogical composition of the target sample; ii) the characteristics and geometry of the emitting X-ray source and detector relative to the sample; iii) the detector's specs and material.
The generated spectra are then convolved to match the anticipated response profile of the rover's spectrometer.

\section{Model Architecture and Physics-Based Attention}
The core architecture of our framework is composed of a CNN backbone and a regression layer; specifically, the approach works with any convolutional backbone (assuming the last layer is re-trained if swapped)
%like SE Resnet 18 (\citet{Farisi2025})
; for our experiments we used a Sequential CNN with 3 Convolutional blocks (using \texttt{BatchNorm} and \texttt{MaxPool}).

\subsection{Multi-Head Regression and Attention Masks}
Instead of using a single linear layer, our architecture features multiple regression heads, each predicting the concentration of a specific element; this design allows us to inject physical knowledge in the form of \textbf{element-specific attention masks}. Each mask is multiplied element-wise with the activations of the final convolutional layer during each forward pass, increasing focus on relevant parts of the input with positive or negative weights (and suppressing irrelevant background with values close to 0). Without this additional step, unconstrained models suffer from gradient instability and often fail to converge, resulting in negative $R^2$ scores on the Test set.

\subsection{Physics-Based Initialization}
The attention masks are not fixed, rather they are learned at training time; however, they are not initialized randomly as all other weights. 
%Since each element emits photons at a precise energy level, an XRF Analysis spectrum will have peaks near the elements' \textit{emission lines}: 
By initializing the masks to higher values in the intervals containing the relevant XRF emission lines, we can guide Gradient Descent towards convergence. To account for any overlap between elements, negative weights are assigned to parts of the input associated with interfering elements. 

The product of learned mask and feature weights is a heatmap that indicates which intervals of the original input the model relies upon for its estimations, providing a form of explainability for the model.

\section{Instance-Based Explainability via Deep k-NN}
Our second explainability mechanism is the addition of a separate, interpretable model in the form of a \textbf{Deep K-Nearest Neighbor regressor}; this approach was first proposed by \citet{DeepKNN} in the context of image classification. The algorithm operates on the latent space of the trained CNN and provides instance-based explanations for each query of the black-box model, while also acting as a predictor itself.

Convolutional architectures naturally preserve local spatial features in their representations. After training the backbone network with the regression heads as the last layer, the backbone is frozen and the heads are replaced by the k-NN regressor operating directly on the latent embedding space. The k-NN algorithm predicts concentrations by interpolating the $k$ nearest samples in the training dataset weighted according to a simple Euclidean distance.
To ensure that the latent space is well-ordered (and therefore suitable for k-NN regression) without degrading predictive accuracy, we use a custom loss function combining the Mean Squared Error ($\mathcal{L}_{\text{MSE}}$) of the predictions and a Triple Loss ($\mathcal{L}_{\text{Triplet}}$) on the embeddings:
\begin{equation}
    \mathcal{L}_{\text{total}} = \mathcal{L}_{\text{MSE}} + R \cdot \lambda \cdot \mathcal{L}_{\text{Triplet}}
\end{equation}
where $R$ is a re-scaling factor to be calibrated once per dataset (since the triplet loss depends on the magnitude of the inputs while the MSE doesn't) and $\lambda$ is a balancing hyperparameter between the two, preventing the triplet regularization from overwhelming the primary regression objective.

\section{Experimental Evaluation}
Following is a series of experiments aimed to show the predictive performance of our approach both with and without the k-NN regressor w.r.t. different hyperparameters \& concentrations. All results are the average of 10 complete training runs.

Our first experiment averages the results on multiple metrics for the regression heads and the k-NN algorithm over 10 different trainings on the same dataset. The results in Table 1 show a minimal decrease on the $R^2$ score, going from an average of 0.95 to 0.91; this can be considered negligible, as the k-NN is mainly employed as an explainability mechanism. It's worth noting, however, how it could also be usable as a standalone predictor. 

The second test aims to investigate the impact of the $\lambda$ parameter in the Loss function.
A value for $\lambda$ too low does not order the latent space enough for k-NN to function, while a value too high hinders the regressor’s performance (see Table 2). With a value of 0.2, the drop in performance between the two methods can be considered negligible, as the k-NN is mainly employed as an explainability mechanism. 
The numbers suggest that $\lambda$ isn't simply a trade-off between the two methods, rather it optimizes the two algorithms at the same time and has a sweetspot for optimal performance for both.

\begin{table}[!ht]
\label{table1}
\begin{center}
{\small
\begin{tabular}{|c|c|c|}  % l = left, c = centered
\tableline
Regressor & Average $R^2$ Score & Average RMSE\\
\tableline
Multi-Heads & $0.9406$ & $0.007738$ \\
Deep k-NN & $0.9101$ & $0.017008$ \\
\tableline
\end{tabular}
\caption{Note how the k-NN's RMSE is more than double that of the Heads; this suggests that while both models correctly follow the trends of the data, with k-NN individual errors can still be noticeable and exact quantities are not always identified.}
}
\end{center}
\end{table}
\vspace{-0.5cm}
\begin{table}[!ht]
    \begin{center}
    {\small
    \label{table2}
    \begin{tabular}{|c|c|c|}
        \hline
        $\lambda$ Parameter & Regression Heads & k-NN \\
        \hline
        0.05 & 0.9443 & 0.8796 \\
        0.1  & \textbf{0.9538} & 0.8964 \\
        0.2  & 0.9506 & \textbf{0.9100} \\
        0.3  & 0.9058 & 0.9080 \\
        \hline
    \end{tabular}}
    \caption{Average Performance Comparison across $\lambda$ Values}
\end{center}
\end{table}

\vspace{-0.5cm}


Our final test focused on elements present in extremely small percentages inside the samples: the model kept a consistent performance for quantities down to $0.01\%$, but below that threshold the $R^2$ score quickly fell below 0 (meaning predictions were worse than guessing the average). Lab experiments show that elements in concentrations between $10^{-4}$ and $10^{-5}$ like Thorium (typically around 15 parts per million) were not detectable by our instrument, meaning that the behavior of our simulation is accurately representing a real scenario. 

\section{Conclusion and Future Work}
In this work, we presented a framework for spectral regression made explainable by CNN latent space k-NN regressor and physics-informed attention masks. Our approach achieves high predictive accuracy (with an $R^2$ score around 0.9) while providing both model explanations and instance-based interpretability. Future work will focus on applying TinyML principles to run the algorithm on edge computing architectures, improving the quality and realism of the data simulation pipeline with additional noise factors and extensive testing for ensuring robustness of the approach.


% These lines show examples of subject index entries. At this stage these have to commented
% out, and need to be on separate lines. Eventually, they will be automatically uncommented
% and used to generate entries in the Subject Index at the end of the Proceedings volume.
% Don't leave these in! - replace them with ones relevant to your paper.
%\ssindex{conference!ADASS 2026}
%\ssindex{organisations!ASP}


\bibliography{P216}  % For BibTex

\end{document}
